\chapter{数据，问题的载体：数据思维入门}

\section{数据是什么：从现象到数字的映射}

数据是现实世界现象经过观测、定义和记录之后形成的表示。交通领域常见数据包括固定检测器数据、GPS 轨迹、问卷调查数据、遥感数据以及网约车或地图 API 数据。

每类数据都有由采集机制带来的内生缺陷，例如检测器会漂移、GPS 会受到遮挡、调查数据存在偏差。

\section{描述你的数据：告诉AI数据长什么样}

\subsection{数据从哪里来}

采集机制决定数据含义。检测器数据、GPS 轨迹、订单日志、刷卡数据、问卷调查、事故记录和道路视频都可能包含不同类型的偏差。

\subsection{一行数据代表什么}

一行数据可以代表一辆车、一笔订单、一个路段时间片、一个司机日或一个站点小时。观测单位不同，模型问题也完全不同。

\subsection{标签是如何生成的}

标签并非天然正确。事故严重程度可能受记录规范影响，司机接单标签可能受到派单机制影响，拥堵标签则可能受到阈值选择影响。

向 AI 描述数据时，应至少说明数据来源、行和列的语义、时空范围、采样频率、缺失情况和标签含义。

\section{数据的质量问题：AI无法替你解决的部分}

\begin{itemize}
  \item 缺失值背后是什么：传感器故障还是真实零流量。
  \item 异常值是错误还是信号：事故造成的流量骤降还是传感器漂移。
  \item 时间是否正确对齐：不同数据来源可能使用不同的时区和采样频率。
\end{itemize}

\section{与AI协作进行探索性数据分析}

可以让 AI 生成探索性数据分析代码并协助解读结果，但需要通过 Prompt 引导 AI 关注真正重要的业务问题，并人工检查生成代码和结论。

\section{实战练习：数据集的EDA全流程}

选择一个真实数据集，完整描述其来源、观测单位、字段语义和质量问题，再与 AI 协作完成探索性数据分析。
